\par\noindent\textbf{(1)} 交换两出射相同费米子等价于把质心系散射角 \(\theta\) 变为 \(\pi-\theta\)。对于自旋对称的同向自旋三重态，空间波函数必须反对称，因而物理振幅为 \(f_F(\theta)=f(\theta)-f(\pi-\theta)\)，不能直接用可辨粒子的单一散射幅。
\par\noindent\textbf{(2)} 微分截面与振幅模平方成正比，\(|f_F|^2=|f(\theta)|^2+|f(\pi-\theta)|^2-2\operatorname{Re}[f(\theta)f^*(\pi-\theta)]\)。最后一项是量子交换干涉，具有相位依赖，不能简单地把可辨粒子截面乘2。
\par\noindent\textbf{(3)} 在 \(\theta=\pi/2\) 时两路径角度相同，所以 \(f_F(\pi/2)=f(\pi/2)-f(\pi/2)=0\)。这一节点属于完全同向自旋费米子和所假定弹性单通道中心势的交换对称性；不同自旋组成的散射不一定有同一节点。
\par\noindent\textbf{(4)} 对对称空间态的同种玻色子或反对称自旋单态费米子，空间振幅是 \(f_B(\theta)=f(\theta)+f(\pi-\theta)\)，90度处增强。对于未极化两自旋 \(1/2\) 费米子，必须按自旋单态、三重态权重分别求截面，且总截面积分要注意相同末态的重复计数。
